% =====================================================================
%  The N-Frame Drag and a Conditional Super-Linear Lower Bound:
%  a formalized reduction to a KRW-style composition hypothesis.
%
%  Self-contained section.  Drop after the Route~C / N-frame drag material.
%  Every lemma tagged [Lean: <module>.<name>] is machine-checked in the
%  repository with `#print axioms` reporting exactly
%      [propext, Classical.choice, Quot.sound]
%  (i.e. no `sorry`, no extra axioms).  Modules live under
%      PallLean/Paper93/DeepMath/PathB/ComputationalDepthNFrame*.lean
%
%  HONEST SCOPE.  The main theorem is CONDITIONAL on Assumption (KRW-C).
%  Everything else is proved unconditionally in Lean.  Even granting
%  (KRW-C), the conclusion is a super-linear lower bound on `cbudget`
%  (gate count) for one explicit family in the CGate model.  It is NOT
%  `NEXP \not\subset ACC^0` and NOT `P \neq NP`.
% =====================================================================

\section{The N-frame drag: a linear method, and its conditional escape}
\label{sec:nframe-krw-reduction}

We record, in one place, the exact reach of the N-frame drag and the single
hypothesis under which it yields a super-linear bound.  The structure is a
\emph{reduction}: we prove, unconditionally and in Lean, a chain of lemmas that
isolates one open statement---a composition non-collapse hypothesis in the
Karchmer--Raz--Wigderson (KRW) family---and we show that this hypothesis is
\emph{sufficient} to drive the drag past its linear ceiling to $\Theta(N\log N)$.

Throughout, $\mathrm{cbudget}(f)$ is the minimal CGate count of $f$;
$K=|\mathrm{ESS}(f)|$ its essential-variable count; and
$\mathrm{coneExcess}(C)=\sum_{w}(\mathrm{readers}(w)-1)$ the total excess fan-out
of a circuit $C$.  The N-frame ledger (\emph{connectivity--fanout}) states, for
every circuit $C$ computing $f$,
\begin{equation}
  2K + \mathrm{coneExcess}(C) \;\le\; \mathrm{length}(C)+1 ,
  \tag{Ledger}
  \label{eq:ledger}
\end{equation}
so a lower bound on $\mathrm{coneExcess}$ for \emph{every} circuit lifts to a
lower bound on $\mathrm{cbudget}$.

\subsection{The linear ceiling}

The drag's only handle on $\mathrm{coneExcess}$ is \emph{cut capacity}: at a
balanced wire cut, a family $Y\subseteq\{0,1\}^N$ of pairwise-distinguished rows
forces $|Y|\le 2^{\,\mathrm{coneExcess}+1}$, i.e.\ $\mathrm{coneExcess}\ge
\log_2|Y|-1$.  This is capped.

\begin{lemma}[Row-family bound; linear ceiling]
\label{lem:drag-ceiling}
Any distinguished row family $Y\subseteq\{0,1\}^N$ has $|Y|\le 2^N$, hence
$\log_2|Y|\le N$.  Consequently the drag-provable bound satisfies
$2K+\mathrm{coneExcess}\le 3N$ whenever $K\le N$ and the cut-capacity certificate
$\mathrm{coneExcess}\le N$.
\end{lemma}
\noindent\emph{[Lean: \texttt{NFrameDragCeiling.rowFamily\_card\_le},
\texttt{NFrameDragCeiling.drag\_linear\_ceiling}.]}

Thus the N-frame drag is \emph{structurally a linear method}: it proves
$\mathrm{cbudget}\ge(2+c)N$ and provably cannot exceed $\sim 3N$ through a single
cut.  Any super-linear conclusion must come from a mechanism that is not a single
cut-capacity certificate.

\subsection{The escape is a recursion, not a bigger cut}

A single cut is capped at $N$.  A \emph{recursion} accumulates many sub-$N$
certificates across $\log N$ scales without any single cut exceeding its level's
input count.

\begin{lemma}[Amplification arithmetic]
\label{lem:amplify}
Let $T:\mathbb N\to\mathbb N$ and $c\ge 1$.  If
$2\,T(k)+c\cdot 2^{\,k+1}\le T(k+1)$ for all $k$, then $c\,(k\,2^{k})\le T(k)$ for
all $k$.  In particular, writing $N=2^{k}$, one has $T(k)\ge c\,N\log_2 N$, which
exceeds every linear bound $b\cdot N$.
\end{lemma}
\noindent\emph{[Lean: \texttt{NFrameConeAmplify.coneExcess\_amplify},
\texttt{NFrameConeAmplify.amplify\_exceeds\_linear}.]}

So \emph{if} the coneExcess $T(k)$ of a recursive family obeys the recurrence
$T(k{+}1)\ge 2\,T(k)+c\cdot 2^{k+1}$, the drag yields $\Theta(N\log N)$.  The rest
of this section establishes every ingredient of that recurrence except one, which
becomes Assumption~(KRW-C).

\subsection{The recursive family}

Fix a constant-degree Ramanujan graph and the associated quadratic detection form
$q(x)=x^\top A x$ over $\mathbb F_2$ with polarization
$\mathrm{bilinSym}(A;x,\delta)=x^\top(A{+}A^\top)\delta$; write $M=A+A^\top$.  The
family is defined by recursion on $k$ (size $N=2^{k}$):
\[
  f_{k+1}(x) \;=\; g_{k}\!\bigl(f_{k}(x_{1}),\,f_{k}(x_{2})\bigr),
  \qquad x=(x_{1},x_{2}),
\]
with $x_1,x_2$ disjoint halves, $f_k:\{0,1\}^{N}\to\{0,1\}^{M}$, and $g_k$ a
\emph{multi-output injective mixing} (Subsection~\ref{ss:firewall}); $f_0$ is the
base gadget below.

\subsection{The base case: two disjoint rigidities add, and cannot cancel}

At the base level the family is the expander quadratic form, and the direct sum
holds \emph{unconditionally}.

\begin{lemma}[Detection direct sum]
\label{lem:direct-sum}
If $s_1,t_1$ and $s_2,t_2$ are induced matchings on disjoint blocks---each with
identity detection $\mathrm{bilinSym}(A;e_{t_i k},e_{s_i l})=[k{=}l]$ and zero
cross-detection---then the combined matching over $\mathrm{Fin}\,r_1\oplus
\mathrm{Fin}\,r_2$ has identity detection.  Hence the direct sum is pairwise
distinguished by $q$ over $\mathrm{Fin}(r_1{+}r_2)$: an $(r_1{+}r_2)$-fold identity
submatrix of $M$, so $\mathrm{rank}_{\mathbb F_2}(M_1\oplus M_2)\ge r_1+r_2$.
\end{lemma}
\noindent\emph{[Lean: \texttt{NFrameDirectSum.bilinSym\_cross\_zero},
\texttt{NFrameDirectSum.combined\_detection\_identity},
\texttt{NFrameCancellation.direct\_sum\_distinct},
\texttt{NFrameInducedMatch.induced\_matching\_distinct}.]}

\begin{lemma}[Cancellation is ruled out for $q$]
\label{lem:cancellation}
$\mathbb F_2$ cancellation is a linear change of variables, under which
$\mathrm{rank}_{\mathbb F_2}(M)$ is invariant; and $\mathrm{rank}$ is additive
under direct sum (Lemma~\ref{lem:direct-sum}).  By Mirwald--Schnorr the
multiplicative complexity of a quadratic form over $\mathbb F_2$ is
$\ge \mathrm{rank}(M)/2$.  Hence no linear mixing of two disjoint copies of $q$
reduces the AND-gate count below the sum; the base-case direct sum is
cancellation-proof.  Structurally, disjoint copies decompose additively:
$q(x{+}y)=q(x)+q(y)$ whenever $\mathrm{bilinSym}(A;x,y)=0$.
\end{lemma}
\noindent\emph{[Lean: \texttt{NFrameUhligCheck.qform\_additive\_disjoint},
\texttt{NFrameUhligCheck.qform\_additive\_pair}; Mirwald--Schnorr cited.]}

\begin{remark}[Uhlig does not bind]
Uhlig mass production reduces the direct sum only for near-maximally-hard
functions ($\mathrm{cbudget}\sim 2^n/n$).  Here $q$ is linear-size and $f_k$
quasi-linear ($O(N\log N)$); neither is in that regime, so Uhlig does not apply.
\end{remark}

\subsection{Why the base case does not lift: the two failure modes}

Lemmas~\ref{lem:direct-sum}--\ref{lem:cancellation} are \emph{rank} statements,
hence capped at $N$.  A super-linear step needs a \emph{non-rank} certificate for
$\mathrm{coneExcess}$.  There is no certificate for $\mathrm{coneExcess}$ in
isolation---every function has a formula, with $\mathrm{coneExcess}=0$---so the
object is the length$\leftrightarrow$coneExcess tradeoff, and the two available
certificates fail in orthogonal ways.

\begin{lemma}[Failure mode I: rank cap]
\label{lem:rank-cap}
Cut capacity gives $\mathrm{coneExcess}\ge \mathrm{cut\text{-}rank}\le
\log_2|Y|\le N$: capped at the input dimension.
\end{lemma}
\noindent\emph{[Lean: \texttt{NFrameDragCeiling.rowFamily\_card\_le}.]}

\begin{lemma}[Failure mode II: formula collapse]
\label{lem:formula-collapse}
A formula lower bound $F$ transfers to circuits only through the unfolding loss
$\mathrm{formulaSize}\le \mathrm{length}\cdot 2^{\,\mathrm{coneExcess}}$, i.e.\
$\mathrm{length}\ge F/2^{\,\mathrm{coneExcess}}$.  At
$\mathrm{coneExcess}=\lceil\log_2 F\rceil$ this bound is $\le 1$, and the best
guaranteed $\max(\mathrm{coneExcess},\,F/2^{\,\mathrm{coneExcess}})$ over the
designer's fan-out choice is $\le \lceil\log_2 F\rceil$.  The formula route
certifies at most $\mathrm{length}\ge\log_2 F$---logarithmic---even from a
super-linear $F$.
\end{lemma}
\noindent\emph{[Lean: \texttt{NFrameNonRankCert.formula\_cert\_collapse},
\texttt{NFrameNonRankCert.formula\_cert\_ceiling}.]}

A working certificate must dodge both: uncapped by input dimension \emph{and}
stable under fan-out.

\subsection{The demand certificate dodges both, at the price of one new requirement}

Let a value be \emph{demanded} at $K$ sites if any circuit for $f$ must make it
present there.  Such a value must be \emph{routed} (fan-out $K{-}1$, charged to
$\mathrm{coneExcess}$) or \emph{recomputed} ($K{-}1$ extra copies at unit cost
$c\ge 1$, charged to $\mathrm{length}$); both land in \eqref{eq:ledger}.

\begin{lemma}[Demand certificate]
\label{lem:demand}
Recomputing never beats routing: $d\le d\cdot c$ for $c\ge 1$, so a demand of $K$
costs $\ge K{-}1$ under either strategy (no $2^{-\mathrm{coneExcess}}$ decay, so it
dodges Failure~II).  No route/recompute mix beats the total: for charges
$\ge d_i$, $\sum_i d_i\le\sum_i \mathrm{charge}_i$.  The demand total is uncapped:
it exceeds every $N$ (dodging Failure~I).  Given irreducibility (each charge
$\ge d_i$) and \eqref{eq:ledger}, $2K+\sum_i d_i\le \mathrm{length}$.
\end{lemma}
\noindent\emph{[Lean: \texttt{NFrameDemandCert.recompute\_charge},
\texttt{NFrameDemandCert.demand\_no\_strategy\_beats\_total},
\texttt{NFrameDemandCert.demand\_uncapped},
\texttt{NFrameDemandCert.demand\_certificate}.]}

The demand certificate dodges both failure modes but introduces one new
requirement---\emph{irreducibility under sharing}: demand can be reduced by
computing shared intermediate values once.  For the linear map $x\mapsto Mx$ this
is exactly the linear-circuit complexity of $M$, and proving it super-linear is
Valiant's rigidity program.  Thus the demand certificate does not evade the open
problem; it relocates it to irreducibility.

\subsection{A single expander cannot supply super-linear demand}

\begin{lemma}[Sparsity caps demand]
\label{lem:expander-cap}
A $d$-regular expander has total demand $\sum_{i} d = d\cdot N$ (each input reused
$\deg=d$ times), which for bounded degree ($d<N$) satisfies $d\cdot N<N^{2}$:
strictly below the super-linear regime.  The associated demand certificate yields
only $2K+d\cdot N\le \mathrm{length}$---linear.
\end{lemma}
\noindent\emph{[Lean: \texttt{NFrameExpanderDemand.regular\_demand\_total},
\texttt{NFrameExpanderDemand.regular\_demand\_below\_superlinear},
\texttt{NFrameExpanderDemand.expander\_demand\_linear}.]}

The sparsity that makes the expander explicit and every-cut rigid is exactly what
caps its demand at $O(N)$.  Hence super-linear demand cannot come from a single
bounded-degree matrix; the only remaining route is \emph{composition}.

\subsection{Composition and the firewall}
\label{ss:firewall}

Fix $x_2=c$, so $f_{k+1}(\cdot,c)=g_k\!\bigl(f_k(\cdot),f_k(c)\bigr)$.  Whether
this restriction forces the sub-instance is governed by injectivity of the mixing.

\begin{lemma}[Firewall condition]
\label{lem:firewall}
If $g_k(\cdot,b)$ is injective, then the restriction distinguishes every pair
$x,x'$ that $f_k$ distinguishes: $f_k(x)\ne f_k(x')\Rightarrow
g_k(f_k(x),b)\ne g_k(f_k(x'),b)$.  Conversely, if $g_k$ collapses two distinct
sub-outputs then $g_k(\cdot,b)$ is not injective and the firewall fails.  Moreover
for $M\ge 2$ every single-output map $(\mathrm{Fin}\,M\to\mathbb F_2)\to\mathbb
F_2$ is non-injective ($2^{M}>2$); so the mixing $g_k$ \emph{must be
multi-output} ($M$-bit), injective in each argument---e.g.\ $a\mapsto a\oplus b$
or $a\mapsto a\oplus Bb$, not a single-output quadratic form.
\end{lemma}
\noindent\emph{[Lean:
\texttt{NFrameMixingFirewall.injective\_mixing\_preserves\_distinct},
\texttt{NFrameMixingFirewall.firewall\_restriction\_distinguishes},
\texttt{NFrameMixingFirewall.noninjective\_mixing\_collapses},
\texttt{NFrameMixingFirewall.single\_output\_mixing\_not\_injective}.]}

An injective multi-output mixing thus \emph{forces each sub-instance}, but only
$1\times$: one restriction re-uses the whole circuit.  Upgrading $1\times$ to
$2\times$---the two forced sub-demands being \emph{disjoint} across scales, not
shared through $g_k$---is precisely a composition-does-not-collapse statement.  We
isolate it.

\begin{assumption}[KRW-style composition non-collapse; (KRW-C)]
\label{ass:krw-c}
For the recursive family $\{f_k\}$ with a fixed multi-output injective mixing
$g_k$ satisfying the firewall condition (Lemma~\ref{lem:firewall}), the
$\mathrm{coneExcess}$ profile $T(k):=\min_{C}\mathrm{coneExcess}(C)$ over circuits
$C$ computing $f_k$ satisfies, for a constant $c\ge 1$ and all $k$,
\[
  T(k+1)\;\ge\;2\,T(k)\;+\;c\cdot 2^{\,k+1}.
\]
Here the $+\,c\cdot 2^{k+1}$ is the fresh mixing demand supplied by expander
rigidity at level $k{+}1$ (Lemma~\ref{lem:expander-cap} gives it at the linear
per-level scale), and the coefficient $2$ is the cross-scale disjointness: no
circuit for $f_{k+1}$ shares $\mathrm{coneExcess}$ between the two $f_k$
sub-instances through $g_k$.
\end{assumption}

\begin{remark}[On (KRW-C)]
The firewall (Lemma~\ref{lem:firewall}) delivers the $1\times$ inclusion
$T(k+1)\ge T(k)+c\cdot 2^{k+1}$ unconditionally; (KRW-C) is exactly the upgrade of
the coefficient from $1$ to $2$.  This is a composition non-collapse in the
$\mathrm{coneExcess}$/demand setting, kin to the Karchmer--Raz--Wigderson
conjecture (which asserts formula depth adds under composition).  We do
\emph{not} claim (KRW-C) is implied by KRW: KRW concerns formula depth ($NC^1$),
whereas (KRW-C) concerns circuit $\mathrm{coneExcess}$ (super-linear size); they
are structurally analogous but formally distinct.  (KRW-C) is open, and proving
it for this explicit family is at least as hard as the corresponding composition
lower bound.
\end{remark}

\subsection{The conditional theorem}

\begin{theorem}[Conditional super-linear drag bound]
\label{thm:conditional-superlinear}
Assume \textup{(KRW-C)} for the recursive family $\{f_k\}$.  Then
$\mathrm{coneExcess}(f_k)\ge c\,k\,2^{k}$, i.e.\ writing $N=2^{k}$,
\[
  \mathrm{coneExcess}(f_k)\;\ge\; c\,N\log_2 N,
  \qquad
  \mathrm{cbudget}(f_k)\;\ge\;2K + c\,N\log_2 N .
\]
In particular $\mathrm{cbudget}(f_k)$ is super-linear in $N$: it exceeds every
linear bound $b\cdot N$, and hence the $3N$ ceiling of Lemma~\ref{lem:drag-ceiling}.
\end{theorem}

\begin{proof}
By (KRW-C) the profile $T(k)=\mathrm{coneExcess}(f_k)$ satisfies the hypothesis of
Lemma~\ref{lem:amplify} with the given $c\ge 1$; hence $T(k)\ge c\,(k\,2^k)=
c\,N\log_2 N$.  The bound on $\mathrm{cbudget}$ follows from \eqref{eq:ledger},
and super-linearity from
\texttt{NFrameConeAmplify.amplify\_exceeds\_linear}.
\end{proof}

\noindent All steps of the proof except the invocation of (KRW-C) are machine
checked; the reduction $\textup{(KRW-C)}\Rightarrow$
Theorem~\ref{thm:conditional-superlinear} is
\texttt{NFrameConeAmplify.coneExcess\_amplify} applied to the assumed recurrence.

\subsection{Honest scope}

Theorem~\ref{thm:conditional-superlinear} is a \emph{conditional} lower bound.
Unconditionally and in Lean we have proved: the drag's linear ceiling
(Lemma~\ref{lem:drag-ceiling}); the amplification arithmetic
(Lemma~\ref{lem:amplify}); the base-case direct sum and its cancellation-proofness
(Lemmas~\ref{lem:direct-sum}--\ref{lem:cancellation}); the two orthogonal failure
modes of every non-demand certificate (Lemmas~\ref{lem:rank-cap}--%
\ref{lem:formula-collapse}); the demand certificate and its two dodges
(Lemma~\ref{lem:demand}); the sparsity cap on single-matrix demand
(Lemma~\ref{lem:expander-cap}); and the firewall condition forcing a multi-output
injective mixing (Lemma~\ref{lem:firewall}).  The one remaining gap---the upgrade
from the proved $1\times$ firewall to the $2\times$ cross-scale disjointness---is
isolated as Assumption~(KRW-C), a composition non-collapse in the KRW family.

We claim nothing beyond this.  The conclusion, even under (KRW-C), is a
super-linear lower bound on $\mathrm{cbudget}$ (CGate count) for one explicit
recursive family.  It is not a separation of complexity classes.  Nothing here is
$\mathrm{NEXP}\not\subset\mathrm{ACC}^0$ or $\mathrm{P}\neq\mathrm{NP}$.
